\section{Dataflow Mapping}
\label{sec:dataflow}

A configuration is only useful if the array can be fed correctly. This section
describes how three classes of computation map onto the fabric, and why the
architectural constraint of Section~\ref{sec:arch} admits some mappings and
excludes others.

\subsection{Element-wise and reduction modes}

The \emph{diagonal} mapping computes element-wise binary operations. Operand
$a$ is injected on the north edge and operand $b$ on the west edge; PEs above
the diagonal pass the $a$ values southward and those below pass the $b$ values
eastward, so that each diagonal PE receives one $(a,b)$ pair and computes one
output. Four elements are produced per transaction, one per diagonal PE.

The \emph{reduction} mapping accumulates a stream into a single PE using the
multiply-accumulate operation: at each step the host injects one pair on the
north and west ports of PE$(0,0)$, which folds the product into its register.
After the stream is exhausted the register holds the dot product.

\subsection{Weight-stationary systolic matrix--vector}

A matrix--vector product $y = A x$ uses all sixteen PEs. Because interior PEs
cannot receive host data directly, the vector $x$ is held \emph{stationary} in
the PE immediates while the matrix is \emph{streamed}: row~0 multiplies each
incoming matrix element by its column weight, and rows~1--3 form a delay line
that shifts the partial products southward. Figure~\ref{fig:systolic} shows the
resulting space--time skew: after $R=\text{ROWS}$ steps, PE row $R{-}1{-}i$
holds the products for output row $i$, which the host sums. Larger matrices are
tiled over $4\times4$ blocks.

\begin{figure}[t]
  \centering
  \resizebox{\columnwidth}{!}{% figures/systolic.tex — weight-stationary systolic matrix-vector (space-time).
\begin{tikzpicture}[x=13mm, y=11mm, font=\scriptsize]
  % left: the array column behaviour (row roles)
  \node[pe, fill=cAlu, draw=cAluE] (r0) at (0,0) {MUL};
  \node[pe] (r1) at (0,-1) {PASS};
  \node[pe] (r2) at (0,-2) {PASS};
  \node[pe] (r3) at (0,-3) {PASS};
  \node[lbl, left=1mm of r0] {row 0};
  \node[lbl, left=1mm of r3] {row 3};
  \draw[wire] (r0)--(r1); \draw[wire] (r1)--(r2); \draw[wire] (r2)--(r3);
  \draw[wire] (0,1) node[lbl,above]{$A[k]$ (streamed)} -- (r0);
  \node[lbl, right=1mm of r0] {$\times\,x_c$ (stationary)};

  % right: space-time skew table. cell = which matrix row's product is held
  \begin{scope}[xshift=42mm]
    \matrix (m) [matrix of nodes, nodes={draw=cWire!40, minimum size=8mm,
                 anchor=center, font=\scriptsize}, column sep=-\pgflinewidth,
                 row sep=-\pgflinewidth]
    {
      $A_3x$ & $A_2x$ & $A_1x$ & $A_0x$ \\
      $A_2x$ & $A_1x$ & $A_0x$ & \\
      $A_1x$ & $A_0x$ &        & \\
      $A_0x$ &        &        & \\
    };
    \node[lbl, above=1mm of m] {step $\rightarrow$};
    \node[lbl, rotate=90, left=1mm of m] {PE row $\downarrow$};
    % highlight the read-out diagonal (after ROWS steps)
    \draw[cRegE, thick, rounded corners]
       (m-1-1.north west) rectangle (m-1-1.south east);
    \node[lbl, right=10mm of m, align=left, cRegE] {after $R$ steps:\\
       $y_i=\sum_c \reg{PE}(R{-}1{-}i,c)$};
  \end{scope}
\end{tikzpicture}
}
  \caption{Weight-stationary systolic matrix--vector product. The matrix
  streams in one row per step; products march south through the PASS delay
  line. The skewed diagonal read after $R$ steps gives each output row.}
  \label{fig:systolic}
\end{figure}

\subsection{Why weight-stationary}

A classical output-stationary systolic array (as in a tensor processing
unit~\cite{tpu,kung}) fuses, in each PE, a multiply with the addition of a
partial sum flowing from a neighbour, while a second stream of activations flows
orthogonally. That requires each PE to forward one value \emph{and} retain a
distinct accumulator---two independent live values. The present fabric provides
a single register per PE, and external data enters only at two edges. Fused
spatial accumulation is therefore not realisable here, whereas the
weight-stationary streaming dataflow---one operand fixed in the immediates, a
single stream marching through PASS cells---is the natural fit. Convolution
follows directly, as it is a matrix--vector product with a Toeplitz matrix; an
inclusive prefix scan is expressed similarly on a single row with the running
value flowing west to east.
